\section*{第 \ref{ch:g6:symmetry} 章 --- 轴对称}

\begin{solution}{exo:g6:symmetry:1}
$A'$ 在 $d$ 的 \emph{右侧}两格，与
$A$ 同行；$B'$ 位于右侧五格、上方一格——
恰为镜像位置。
\end{solution}

\begin{solution}{exo:g6:symmetry:2}
作 $d$ 的垂线，经过 $M$，垂足为 $H$；则 $M'$ 位于
另一侧，且 $HM' = MH = 3$ cm。$MM'$ 的长为
$3 + 3 = 6$ cm。
\end{solution}

\begin{solution}{exo:g6:symmetry:3}
竖直对称轴（同 A 类）：A、H、I、M、O、T、U、V、W、X、Y。
水平对称轴（同 B 类）：B、C、D、E、H、I、K、O、X。
两者兼有：H、I、O、X。
（精确答案视字母的画法而略有不同。）
\end{solution}

\begin{solution}{exo:g6:symmetry:4}
等边三角形有 $3$ 条对称轴；长方形有 $2$ 条（即
两组对边的垂直平分线）；菱形有 $2$ 条（即两条对角线）；正方形有 $4$ 条。
圆有无穷多条对称轴（经过圆心的每一条直线）。
\end{solution}

\begin{solution}{exo:g6:symmetry:5}
\emph{错误。}沿对角线折叠一个（非正方形的）长方形，
两半并不能重合：两个三角形形状相同，
但摆放位置不同——可用一张纸试一试。只有
正方形的对角线才能作为对称轴。
\end{solution}

\begin{solution}{exo:g6:symmetry:6}
反射保持长度、角和周长
（见\cref{prop:g6:symmetry:preserved}）：$A'B' = 4$ cm，
$\widehat{B'A'C'} = 70^\circ$，周长 $A'B'C' = 12$ cm。
\end{solution}

\begin{solution}{exo:g6:symmetry:7}
用圆规作图，方法见\cref{met:g6:symmetry:bisector}。对于直线上的任一
点 $M$，由直尺可验证 $MA = MB$
（见\cref{prop:g6:symmetry:bisector}）。
\end{solution}

\begin{solution}{exo:g6:symmetry:8}
到 $A$ 与 $B$ 距离相等的所有点构成
$[AB]$ 的垂直平分线（见\cref{prop:g6:symmetry:bisector}）。
喷泉可建在该直线上的任一处（花园范围内）。
\end{solution}

\begin{solution}{exo:g6:symmetry:9}
每个顶点都沿垂直于对称轴的方向作反射：对于
$45^\circ$ 的对称轴，位于轴右侧 $k$ 格的点，落在
轴上方 $k$ 格处（反之亦然）——水平与竖直方向的
位移互换。逐一作出各顶点的对称点，然后
按顺序连接。
\end{solution}

\begin{solution}{exo:g6:symmetry:10}
$M$ 的对称点是圆与垂直于
$d$ 且过 $M$ 的直线的第二个交点，即它位于圆上。理由：
反射保持距离，并固定 $O$（位于 $d$ 上），因此
设 $r$ 为任一点到 $O$ 的距离，
则其像点仍在半径为 $r$、圆心为 $O$ 的圆上，即圆被映射到自身。
\end{solution}

\begin{solution}{exo:g6:symmetry:11}
每个交点 $P$ 均由圆规以相同张角 $\rho$
分别从 $A$、$B$ 量得，故 $PA = \rho = PB$。由
\cref{prop:g6:symmetry:bisector}可知，到 $A$、$B$ 距离相等的点
位于 $[AB]$ 的垂直平分线上；两个这样的点
即可确定这条直线。
\end{solution}

\begin{solution}{exo:g6:symmetry:12}
对墙上任一反弹点 $P$，有 $PB = PB'$（关于
$d$ 的反射保持距离），故路径长为
$AP + PB = AP + PB'$。折线 $A \to P \to B'$ 长度最小，
当其为直线时，即 $P$ 位于线段 $[AB']$ 上。故
最佳反弹点为线段 $[AB']$ 与墙的交点。
\end{solution}

\begin{solution}{pb:g6:symmetry:1}
\textbf{1.} 作图方法见\cref{met:g6:symmetry:bisector}；
对直线上任一点 $M$，由直尺可验证 $MA = MB$
（见\cref{prop:g6:symmetry:bisector}）。

\textbf{2.} 喷泉问题需要用到另一个方向的结论：
\emph{所有}与两棵树等距的点都在垂直平分线上，因此
可选位置恰为该直线，别无他处。而
不在垂直平分线上的点必然\emph{不}与两棵树等距：
它严格地更靠近其中一棵。

\textbf{3.} 两条路径共用边 $[MP]$，且 $PB = PA$
（因 $P$ 在垂直平分线上）。故路径 $M \to P \to B$ 与
$M \to P \to A$ 的路径长度相同，均为 $MP + PA$。
现在 $MB$ 是到 $B$ 的直线路线，这里它\emph{恰好}就是
经过 $P$ 的路径，故 $MB = MP + PA$；而到
$A$ 的直线路线比绕经 $P$ 的路线更短：$MA < MP + PA = MB$。
垂直平分线上位于 $A$ 一侧的每一点都更靠近 $A$。

\textbf{4.} 分界线是 $[AB]$ 的垂直平分线：
学校 $A$ 恰好服务于镇上位于该直线一侧的那一半
（第 3 题），学校 $B$ 服务另一半，
位于该直线上的孩子可任选其一。

\textbf{5.} $O$ 位于 $[AB]$ 的垂直平分线上，故 $OA = OB$；
且位于 $[BC]$ 的垂直平分线上，故 $OB = OC$
（见\cref{prop:g6:symmetry:bisector}）。

\textbf{6.} 由 $OA = OB$ 与 $OB = OC$ 得 $OA = OC$。由
\cref{prop:g6:symmetry:bisector}的第二个方向，到 $A$、$C$ 距离相等的点
位于 $[AC]$ 的垂直平分线上，
即第三条垂直平分线不必另问便过 $O$：
任意三角形的三条垂直平分线
\emph{共点}。

\textbf{7.} 以 $O$ 为圆心、$OA$ 为半径的圆，
由全部距离为 $OA$ 的点组成，这些距离皆自 $O$ 量起
（\cref{def:g6:lines:circle}）；由 $OB = OA$、$OC = OA$ 可知，
点 $B$、$C$ 也在圆上。一个圆心、一个半径，
便同时满足三点：它就是经过 $A$、$B$、$C$ 的圆。

\textbf{8.} 两条垂直平分线即已足够，因为第三条自动成立
（第 6 题）：它们的交点已满足全部三个距离相等的条件。
圆以 $O$ 为圆心画出，圆规自 $O$ 张开至三个顶点中
的任一顶点即可。

\textbf{9.} 作法：在留存的弧上标出三点 $A$、$B$、$C$；
作弦 $[AB]$ 与 $[BC]$ 的垂直平分线，其交点即为
圆盘的圆心，该点到 $A$ 的距离即为半径。
此法可行，是因为圆盘的边缘是经过这三个标记点的圆，
而第 8 题确定了经过它们的\emph{唯一}圆。
测试弧可
据此完全复原。

\textbf{10.} 若 $A$、$B$、$C$ 共线，则 $[AB]$ 与 $[BC]$ 的
垂直平分线都垂直于同一直线 $(AB) = (BC)$ --- 故二者
互相平行
（\cref{prop:g6:lines:facts}），且二者互异，从而永不相交。
不存在到三点等距的点，
也没有圆经过三个共线的点。

\textbf{11.} 正五边形有 $5$ 条对称轴，每条过一个顶点
及其对边中点（奇数条：每条轴都属于这一种）。
正六边形有 $6$ 条对称轴 --- 三条过相对顶点，
三条过相对边的中点
（偶数条：两类各占一半）。规律是：正
边数为 $n$ 的正多边形有 $n$ 条对称轴。

\textbf{12.} 剪口同时穿过两层，因此两层恰好失去相配的部分；
展开即为沿折痕的反射，它把每一层的边缘
映到另一层的边缘上。孔洞自成镜像：
折痕即是它的对称轴 ---
无论剪掉的形状是什么。

\textbf{13.} 展开后的孔洞至少有 $2$ 条对称轴，即两条折痕，
剪口共出现 $4$ 份
（每展开一次图形便加倍：$1 \to 2 \to 4$）。若像雪花那样
再多次折叠，副本与
对称轴都会继续增加。

\textbf{14.} $F''$ 与 $F$ 的\emph{方向}相同（两次镜面翻转
相互抵消了镜像），大小与高度也相同 --- 它只是把 $F$ 向右平移。
平移距离为 $8$ 个方格，恰为两条轴之间距离的\emph{两倍}。
两面平行的镜子合用时，
根本不产生镜像，
而是\emph{平移}。

\textbf{15.} $F''$ 大小不变，但此时它
\emph{倒置}，是绕两条轴的交点旋转半圈所得：
$F''$ 上的每一点，
与 $F$ 上的对应点关于该中心互为直径两端。
两条互相垂直的镜子复合成\emph{半转} ---
即\cref{ch:g7:central}中的中心对称。
\end{solution}
